\section*{الفصل \ref{ch:g6:plants-are-producers} --- النباتات منتِجة}

\begin{solution}{exo:g6:plants-are-producers:1}
\omterm{def:g6:plants-are-producers:matter}{المادة المعدنية}: مادة العالم غير الحي --- الماء، والصخر (وكذلك
غازات الهواء، \omterm{prop:g4:what-plants-need:needs}{والأملاح المعدنية}). \omterm{def:g6:plants-are-producers:matter}{والمادة العضوية}: مادة تصنعها
\omterm{def:g6:exploring-our-environment:components}{الكائنات الحية} وتتألف منها أجسامها --- الخشب، والعضل (وكذلك السكر،
والصوف، والشحم).
\end{solution}

\begin{solution}{exo:g6:plants-are-producers:2}
المدخلات: ماء \omterm{prop:g4:what-plants-need:needs}{وأملاح معدنية} من التربة، وغاز (ثاني أكسيد الكربون)
من الهواء. \omterm{prop:g3:balanced-diet:jobs}{والطاقة}: \omterm{def:g6:exploring-our-environment:conditions}{الضوء}، تلتقطه المادة الخضراء. والمخرج: \omterm{def:g6:plants-are-producers:matter}{مادة
عضوية} --- سكر، ثم خشب وورق \omterm{def:g1:plants-around-us:parts}{وجذر} وبذر. والمشغل: \omterm{def:g1:plants-around-us:parts}{الورقة} الخضراء،
نهارًا.
\end{solution}

\begin{solution}{exo:g6:plants-are-producers:3}
معدني: الندى، والغرانيت، وملح البحر. وعضوي: شمع النحل، وزيت
الطعام --- \omterm{def:g1:animals-around-us:covering}{وصدفة} الحلزون من صنع الحلزون، مبنية من \omterm{def:g6:plants-are-producers:matter}{مادة معدنية}
رصّها \omterm{def:g1:animals-around-us:animal}{الحيوان}: أي صنعها \omterm{def:g6:exploring-our-environment:components}{كائن حي} من موادّ معدنية. (ويُقبل أي من
الجوابين مع هذا التعليل.)
\end{solution}

\begin{solution}{exo:g6:plants-are-producers:4}
من الهواء والماء في أكثرها --- الغاز المأخوذ من الهواء والماء
الذي شربته \omterm{def:g1:plants-around-us:parts}{الجذور} --- مركَّبين في \omterm{def:g1:plants-around-us:parts}{الأوراق} \omterm{prop:g3:balanced-diet:jobs}{بطاقة} \omterm{def:g6:exploring-our-environment:conditions}{الضوء}؛ ولم تعطِ
التربة إلا مقادير صغيرة من \omterm{prop:g4:what-plants-need:needs}{الأملاح المعدنية}.
\end{solution}

\begin{solution}{exo:g6:plants-are-producers:5}
تأخذ ثاني أكسيد الكربون وتطلق الأكسجين. وكل \omterm{def:g6:exploring-our-environment:components}{كائن حي} متنفس يزفر
الأول ويحتاج إلى الثاني --- ونحن منهم.
\end{solution}

\begin{solution}{exo:g6:plants-are-producers:6}
\omterm{ex:g4:plants-without-seeds:bulbtuber}{درنة} البطاطا \omterm{ex:g4:plants-without-seeds:bulbtuber}{وبصلة} الخزامى (مخزنان تحت الأرض)، ونصفا الفولة
السميكان (\omterm{def:g3:animals-through-seasons:active}{مخزن غذاء} \omterm{def:g2:from-seed-to-plant:seed}{البذرة})، \omterm{def:g1:plants-around-us:parts}{وجذر} الجزرة السمين، والشوندر الملآن
سكرًا --- وكلها احتياطيات عضوية مودَعة من صناعة الصيف.
\end{solution}

\begin{solution}{exo:g6:plants-are-producers:7}
لأن الصناعة تجري على \omterm{def:g6:exploring-our-environment:conditions}{الضوء}: فمن غير ضوء لا \omterm{prop:g3:balanced-diet:jobs}{طاقة} لبناء \omterm{def:g6:plants-are-producers:matter}{المادة
العضوية}، مهما تنتظر من معادن في التربة. \omterm{def:g6:exploring-our-environment:conditions}{والضوء} هو المدخل الوحيد
الذي لا بديل له --- والأملاح توابل لا وقود.
\end{solution}

\begin{solution}{exo:g6:plants-are-producers:8}
لأن الحلقة الأولى في السلسلة يجب أن تدخل \omterm{def:g6:plants-are-producers:matter}{مادة عضوية} جديدة إلى
العالم الحي، \omterm{def:g1:animals-around-us:animal}{والحيوانات} مستهلِكة: فهي لا تستطيع إلا أن تحوّل \omterm{def:g6:plants-are-producers:matter}{مادة
عضوية} صنعها غيرها. \omterm{def:g5:food-webs:roles}{والمنتِجون} وحدهم يصنعونها من موادّ معدنية ---
فالحلقة الأولى خضراء دائمًا.
\end{solution}

\begin{solution}{exo:g6:plants-are-producers:9}
الخبز: القمح --- \omterm{def:g2:from-seed-to-plant:seed}{بذور} \omterm{def:g1:plants-around-us:plant}{نبتة} القمح؛ وينتهي الأثر عند \omterm{def:g1:plants-around-us:parts}{ورقة} القمح.
والزبدة: القشدة --- \omterm{def:g2:animals-grow-up:born}{الحليب} --- البقرة --- العشب؛ وينتهي الأثر عند
\omterm{def:g1:plants-around-us:parts}{ورقة} العشب. والجبن: \omterm{def:g2:animals-grow-up:born}{الحليب} --- البقرة --- العشب من جديد. ثلاثة
أطعمة، وثلاث \omterm{def:g1:plants-around-us:parts}{أوراق} (أو اثنتان)، كلها في \omterm{def:g6:exploring-our-environment:conditions}{الضوء}.
\end{solution}

\begin{solution}{exo:g6:plants-are-producers:10}
الدريس غذاء البقرة --- \omterm{def:g6:plants-are-producers:matter}{مادة عضوية} تحوّلها إلى جسمها. أما السماد
فليس غذاء \omterm{def:g1:plants-around-us:plant}{النبات}: بل يعيد تموين الأملاح \emph{المعدنية} في التربة
وهو يتعفن. ثم يصنع \omterm{def:g1:plants-around-us:plant}{النبات} مادته العضوية بنفسه من تلك المعادن
والهواء والماء \omterm{def:g6:exploring-our-environment:conditions}{والضوء}. فالبقرة تأكل وجبات؛ \omterm{def:g1:plants-around-us:plant}{والنبات} يصنعها.
\end{solution}

\begin{solution}{exo:g6:plants-are-producers:11}
في \omterm{def:g6:exploring-our-environment:conditions}{الضوء} تطلق صناعة \omterm{def:g1:plants-around-us:plant}{النباتات} النهارية أكسجينًا وتأخذ ثاني أكسيد
الكربون، فتوازن \omterm{def:g4:breathing:movements}{تنفس} \omterm{prop:g3:sorting-living-things:groups}{الأسماك} --- فتبادل الغازات يجري في
الاتجاهين ويدور الحوض. أما في الظلام فتتوقف الصناعة بينما يواصل كل
ساكن (\omterm{def:g1:plants-around-us:plant}{والنباتات} معهم) \omterm{def:g4:breathing:movements}{التنفس}: فينفد الأكسجين، ويتكوّم ثاني أكسيد
الكربون، ويختنق العالم المغلق.
\end{solution}

\begin{solution}{exo:g6:plants-are-producers:12}
الصحيح: أن مادة الشجرة لم تأتِ من التربة --- فوزنه أثبت ذلك إلى
الأبد. والناقص: أن الماء لم يكن إلا مكوّنًا واحدًا من مكوّني
الكتلة؛ ولم يكن في وسعه أن يعرف الغاز المأخوذ من الهواء، وهو الذي
يوفّر كثيرًا من مادة الخشب. والإضافة المعاصرة: ثاني أكسيد الكربون
من الهواء (مع \omterm{prop:g4:what-plants-need:needs}{الأملاح المعدنية} القليلة من التربة) --- ماء وهواء
معًا، يركّبهما \omterm{def:g6:exploring-our-environment:conditions}{الضوء}.
\end{solution}

\begin{solution}{pb:g6:plants-are-producers:1}
\textbf{1.} الطنّان \omterm{def:g6:plants-are-producers:matter}{مادة عضوية} صنعتها \omterm{def:g1:plants-around-us:plant}{النباتات} نفسها من موادّ
معدنية --- الماء، وثاني أكسيد الكربون في الهواء، وأملاح التربة ---
\omterm{def:g6:exploring-our-environment:conditions}{والضوء} طاقتها. فالأكياس تعيد تموين المعادن؛ والهواء والمطر
يوصّلان الكتلة؛ \omterm{def:g1:plants-around-us:parts}{والأوراق} تقوم بالصنع.

\textbf{2.} البطاطا: عضو ادّخار تحت الأرض (\omterm{ex:g4:plants-without-seeds:bulbtuber}{درنة}). والجزرة: \omterm{def:g1:plants-around-us:parts}{جذر}
ادّخار. \omterm{def:g2:from-seed-to-plant:seed}{وبذرة} الفاصولياء: \omterm{def:g3:animals-through-seasons:active}{مخزن غذاء} \omterm{def:g2:from-seed-to-plant:seed}{البذرة}، محزوم \omterm{def:g1:plants-around-us:plant}{للنبتة} التالية.
\omterm{def:g1:plants-around-us:parts}{وورقة} الخس: العضو العامل نفسه --- مشغل التقاط \omterm{def:g6:exploring-our-environment:conditions}{الضوء}، يُؤكل طازجًا.

\textbf{3.} المورّد الحقيقي هو \omterm{def:g6:exploring-our-environment:conditions}{الضوء}: فكل غرام من المحصول بُني
\omterm{prop:g3:balanced-diet:jobs}{بطاقة} \omterm{def:g6:exploring-our-environment:conditions}{الضوء}، نهارًا، في \omterm{def:g1:plants-around-us:parts}{الأوراق}.

\textbf{4.} من أجل \omterm{prop:g4:what-plants-need:needs}{الأملاح المعدنية}: فالمحاصيل تحملها بعيدًا،
ومن غير إعادة تموينها يصير \omterm{def:g1:plants-around-us:plant}{النبات} فقير المعادن كالنعناع في الرمل
--- قليلة في الكتلة، لا غنى عنها في الوظيفة.

\textbf{5.} تجربة الخزانة المظلمة (السجين الشاحب): فالقماش أزال
\omterm{def:g6:exploring-our-environment:conditions}{الضوء}، وأسبوع منه يظهر أثره منذ الآن --- شحوب، وامتداد، وصناعة
ضعيفة.

\textbf{6.} \omterm{ex:g4:plants-without-seeds:bulbtuber}{الدرنات} تملؤها الأعالي: \omterm{def:g1:plants-around-us:parts}{فالأوراق} تصنع السكر نهارًا،
والفائض يُنقل إلى أسفل ويُودَع في \omterm{ex:g4:plants-without-seeds:bulbtuber}{الدرنات} المنتفخة. \omterm{def:g1:plants-around-us:parts}{والأوراق}
العاملة الوارفة تعني فائضًا كبيرًا --- أعالٍ جيدة، \omterm{ex:g4:plants-without-seeds:bulbtuber}{درنات} جيدة.

\textbf{7.} الماء. فالصناعة تحتاج إلى مدخلاتها كلها معًا: \omterm{def:g6:exploring-our-environment:conditions}{والضوء}
يوفّر \omterm{prop:g3:balanced-diet:jobs}{الطاقة} لا المادة --- فإذا جفّت \omterm{def:g1:plants-around-us:parts}{الجذور} لم يبقَ ما يُبنى به،
مهما تكن السماء متوهجة.

\textbf{8.} \omterm{def:g1:plants-around-us:parts}{الجذور} المقلوبة في التربة \omterm{def:g6:plants-are-producers:matter}{مادة عضوية} مردودة إليها:
فهي إذ تتعفن تحرّر معادنها للمحاصيل القادمة --- وهي قاعدة الردّ في
\cref{prop:g5:people-and-nature:lasting}، مدفوعةً مادةً نباتية.
والتبذير هو حرقها.

\textbf{9.} \omterm{def:g1:plants-around-us:plant}{نبتة} البطاطا $\leftarrow$ الخنزير $\leftarrow$
الأسرة. وكل مادتها العضوية صُنعت أول مرة في \omterm{def:g1:plants-around-us:parts}{أوراق} \omterm{def:g1:plants-around-us:plant}{نبتة} البطاطا ---
والخنزير والأسرة لم يفعلا إلا أن حوّلاها.

\textbf{10.} المكوّنات: ماء \omterm{prop:g4:what-plants-need:needs}{وأملاح معدنية} من التربة، وثاني أكسيد
الكربون من الهواء --- مركَّبة مادةَ ورق. ودور الشمس بالضبط: \omterm{prop:g3:balanced-diet:jobs}{طاقة}
لا مادة. فسلطتها ماء وهواء وتراب، \emph{بناها} ضوء الشمس ---
والفضل الشعري مستحق، وقائمة المكوّنات خاطئة.

\textbf{11.} \omterm{def:g5:food-webs:roles}{المحلِّلون} --- حياة التربة من الديدان وحشرات القمل
الخشبي نزولًا إلى كائنات حية أصغر من أن تُرى --- يتغذون على
\omterm{def:g1:plants-around-us:parts}{الأوراق} الميتة ويفكّكون مادتها العضوية عائدةً نحو \omterm{def:g6:plants-are-producers:matter}{مادة معدنية}،
فيعيدون تموين التربة التي ستسحب منها المحاصيل التالية.

\textbf{12.} مثلًا: ``منتِجو الحديقة يبنون \omterm{def:g6:plants-are-producers:matter}{مادة عضوية} --- هي
المحصول --- من \omterm{def:g6:plants-are-producers:matter}{مادة معدنية} من التربة والماء والهواء، مستعملين
\omterm{def:g6:exploring-our-environment:conditions}{الضوء} طاقةً لهم.''
\end{solution}
